\documentclass{exam}

\usepackage[margin=1in]{geometry}

\usepackage{comment}
\usepackage{graphicx}
\usepackage{listings}

\pagestyle{headandfoot}
\firstpageheader{CSCI 221}{}{C++ Intro Assignment}
\runningheader{CSCI 221}{}{C++ Intro Assignment}


\begin{document}
\title{Inheritance, STL, Smart Pointers, and Generics Assignment (Learning)}
\author{CSCI 221: Computer Science Fundamentals 2}
\date{Spring 2026}
\maketitle


This assignment is an opportunity to test your understanding of Inheritance, STL, Smart Pointers, and Generics and receive feedback. 
Point values are assigned so that you can differentiate between large and small mistakes, but this assignment does not affect your grade. 
To submit your assignment, please submit your answers, code, and makefile. 
For full credit, your code should check for invalid input. 

\textbf{Due Date:}
Monday, April 27th at 1:00 pm. 

\begin{questions}

\question[30]
\textbf{Inheritance.}
\begin{parts}
\part[4]
Create a Location class that contains an x coordinate, a y coordinate (both private doubles), and a name (a private heap-allocated string). 
Ensure that your class implements standard functions as well as a print function that prints out the location and name. 

\part[4]
Create a Building class that publicly inherits from Location. 
In addition to a Location's fields, a Building should have a floor area in square feet and a number of floors. 

\part[2]
Ensure that your Building class properly initiates its internal Location when it is constructed. 
Make sure you are calling the constructor of the Location. 

\part[2]
Ensure that your Building class properly deallocates its internal Location when it is deallocated. 
Make sure you are calling the destructor of the Location. 

\part[1]
Can you access the x and y coordinates of a Building inside its methods?

\part[1]
Make the x and y coordinates of a Location protected, rather than private. 
Can you access the x and y coordinates of a Building inside its methods now?

\part[2]
Write a client function (it can be main) that creates a Building. 
Can the client function access a Building's x and y coordinates?

\part[4]
Create an OfficeBuilding class that publicly inherits from Building. 
In addition to a Building's fields, an OfficeBuilding should have a number of cubicles and a number of offices (both private integers). 

\part[2]
Modify your Building class so that it protectedly inherits from Location. 
Can your client function access the getters and setters for a Building's x and y coordinates?

\part[1]
Can you access the getters and setters for an OfficeBuilding's x and y coordinates in its methods?

\part[2]
Modify your Building class so that it privately inherits from Location. 
Can you access the getters and setters for an OfficeBuilding's x and y coordinates in its methods now?

\part[2]
Modify the print method of a Building to also print the building's area and number of floors. 

\part[1]
What does calling print on a Building do?

\part[1]
What does calling print on a *Location that points to a Building do?

\part[1]
Modify your Location function's print method to be virtual. 
Now what does calling print on a *Location that points to a Building do?

%Double inheritance
\end{parts}

%\question[10]
%\textbf{Smart Pointers.}
%For the following, make sure to use C++ style rather than C style. 
%\begin{parts}
%\part[2]
%Create a pointer to an int. 
%Don't initialize it. 
%Write code that checks the value of that pointer correctly. 
%
%\part[2]
%Write code that allocates an int and sets the pointer to point to that int. 
%
%\part[2]
%Write code that deallocates the int you just allocated. 
%
%\part[2]
%Write code that allocates an array of ints of length 17 and sets the pointer to point to that array. 
%
%\part[2]
%Write code that deallocates the int array you just allocated. 
%
%\end{parts}

\question[30]
\textbf{STL and Smart Pointers.}
\begin{parts}
\part[10]
Create a Tracer class with the following functionality:
\begin{enumerate}
\item A Tracer should contain an unsigned integer. 
Our goal is for each Tracer's contained integer to be unique. 

\item A default constructor. 
This should set the Tracer's contained integer to $i$, where $i$ is the number of Tracers that exist at the time $i$ is created. 
You should be able to do this by having a static unsigned integer in the class that is incremented when Tracers are created and decremented when Tracers are destroyed. 
In addition, the default constructor should print out that a Tracer with ID $i$ has been created. 

\item A copy constructor. 
This should similarly create a Tracer with a new unique value using the static integer. 
In addition, the copy constructor should print out that a Tracer with this ID was created by copying the previous Tracer with its ID. 

\item A destructor. 
In addition to managing the static integer, the destructor should print out that a Tracer with ID $i$ has been destroyed. 

\item An = operator. 
The = operator should not modify the stored integer, but it should print out that Tracer with its ID has been set equal to the Tracer with the other ID. 

\item A $<$ operator. 
The $<$ operator should compare the IDs of the Tracers. 
\end{enumerate}

\part[2]
Create an array containing 16 Tracers. 
How many times is each Tracer constructed?

\part[2]
Create an STL vector containing 16 Tracers. 
How many times is each Tracer constructed?

\part[2]
Create an STL set containing 16 Tracers. 
How many times is each Tracer constructed?

\part[3]
Run the STL sort algorithm on your vector of Tracers. 
What happens during the algorithm?

\part[3]
Create another STL set containing 2 Tracers. 
Run the STL set union algorithm on the two sets. 
What happens during the algorithm?

\part[3]
Create an STL set containing 16 unique pointers to Tracers. 
How many times is each Tracer constructed?

\part[3]
Create an STL set containing 16 shared pointers to Tracers. 
How many times is each Tracer constructed?

\part[2]
Why doesn't it make sense to create an STL set containing 16 weak pointers to Tracers?
\end{parts}

\question[10]
\textbf{Templates.}
\begin{parts}
\part[4]
Write a template function called duplicate that returns a pointer to a heap-allocated object that is a copy of the object passed in. 
You should make use of the copy constructor of that object. 

\part[6]
Create a CountedCopies class that contains an object and an unsigned integer that tracks how many copies of that object it is managing. 
(This integer represents an abstract concept. The CountedCopies object should only store one copy of the internal object.)
Your class should contain standard functionality. 
In addition, it should contain the following methods:
\begin{enumerate}
\item A method that returns the number of managed copies available. 
\item A method that adds a number of managed copies. 
\item A method that removes a number of managed copies. 
\item A method that if the internal integer is positive, decrements that integer by one and creates a copy of the internal object and returns a pointer to that copy. 
If the internal integer is 0, it returns nullptr. 
\end{enumerate}

\end{parts}

\end{questions}

\end{document}